\section{Supported deletion in the twenty-vertex block}\label{app:twenty-support}

We give the local case proof used to determine the upper star-forming number
of the cyclic family. In this appendix $F=F_{20}$, with ports $4,12$.
A selected vertex is \emph{active} if it has a selected neighbour in $F$.
For $C\subseteq V(F)$, a vertex $v$ is \emph{supported} if
\begin{enumerate}
\item $v\in C\setminus\{4,12\}$ and $N_F(v)\cap C\ne\varnothing$;
\item every $u\in N_F(v)\cap C$ has a neighbour in $C\setminus\{v\}$;
\item every $x\in N_F(v)\setminus C$ has at least three neighbours in $C$,
 or exactly two, with the other selected neighbour different from $v$ active.
\end{enumerate}
These are the conditions of Lemma~\ref{lem:supported-deletion}.

\begin{lemma}\label{lem:twenty-support}
Every subset of $V(F)$ of size at least thirteen contains a supported vertex.
\end{lemma}
\begin{proof}
Supportedness is upward closed by Lemma~\ref{lem:supported-deletion}, so assume
$|C|=13$. Suppose no supported vertex exists. Write
\[
 H=\{0,1,2,3\},\quad X=\{4,\ldots,11\},\quad
 A=\{12,13,16,19\},\quad B=\{14,15,17,18\}.
\]
The square sides are $\{12,13\}\mid\{16,19\}$ and
$\{14,15\}\mid\{17,18\}$. Let $s(x)$ be the companion witness at the same
helper and $m(x)$ its witness-matching mate:
\[
 s=(4\,5)(6\,7)(8\,9)(10\,11),\qquad
 m=(4\,11)(5\,9)(6\,8)(7\,10).
\]
For $x\ge5$, $N_F(x)=\{h(x),m(x),x+8\}$, where
$h(x)=\lfloor(x-4)/2\rfloor$, whereas $N_F(4)=\{0,11\}$.
Put $h=|C\cap H|$, $k=|C\cap(A\cup B)|$, $D=X\setminus C$ and $p=|D|$.
Counting selected vertices gives
\begin{equation}\label{eq:twenty-budget}
 p=h+k-5.
\end{equation}

\paragraph{Local owner rules.}
Call a nonport owner \emph{ready} if its square is fully selected, or it is
an endpoint of the selected three-vertex path in a square with three
selected vertices. Its square-neighbour tests are automatic. In the
three-vertex case the selected square neighbour retains the other endpoint,
and the omitted square neighbour retains that active endpoint.
A ready owner can therefore fail only at its witness. The possibilities are
\[
\begin{array}{c|c}
\text{square selection}&\text{ready owners}\\\hline
A\text{ full}&13,16,19\\
A\text{ omits }12\text{ or }13&16,19\\
A\text{ omits }16\text{ or }19&13\\
B\text{ full}&14,15,17,18\\
B\text{ omits }14\text{ or }15&17,18\\
B\text{ omits }17\text{ or }18&14,15
\end{array}
\]
A witness-matching pair incident to $t$ ready owners requires at least $t$
omitted witnesses. For $t=1$, selecting both witnesses makes the owner safe.
For $t=2$, if only one witness is omitted, its selected mate is active through
its own selected owner, again making the first owner safe. The pairs are
disjoint, so
\begin{equation}\label{eq:twenty-ready}
 p\ge\#\{\text{ready owners}\}.
\end{equation}
If $x+8$ is ready and exactly one of $x,m(x)$ is omitted, then $h(x)$ must
be omitted. Otherwise a selected $x$ retains its helper, or an omitted $x$
has three selected neighbours. If $x$ is the omitted one, its selected mate
must also be inactive.

When a square has exactly two selected vertices, every omitted owner is a
safe outside neighbour of its selected witness: either both core neighbours
are selected, giving three selected neighbours including that witness, or
one core neighbour is selected and active through the other selected owner.
Also, a selected owner whose same-side opposite is omitted has a selected
core neighbour. We will use these two-owner rules below.

\paragraph{Local helper rule.}
The witness matching gives the helper cycle $0-2-1-3-0$. Represent an
omitted witness by an arrow from its helper to the helper of its mate.
An active selected helper is supported using the twelve-cycle alone if
neither matching mate of its witnesses is omitted and each omitted witness
at that helper has its mate at a selected helper. Indeed, selected witnesses
retain their matching mates, while an omitted witness retains its active
mate after helper deletion. Thus a failed active helper must receive an
omission arrow, or emit one to an omitted helper. An arrow accounts for at
most one active selected helper: its target when selected, otherwise its
source. Consequently
\begin{equation}\label{eq:twenty-helpers}
 \#\{\text{active selected helpers}\}\le p.
\end{equation}
An isolated selected helper requires two omitted witnesses. Hence $h=2,3$
require $p\ge2$, and $h=4$ requires $p\ge3$. If $h=3,p=2$, one selected
helper is isolated and its two arrows target the other selected helpers.
The omitted helper is opposite it, leaving exactly
\begin{equation}\label{eq:twenty-h3-p2}
\begin{array}{c|cccc}
\text{omitted helper}&0&1&2&3\\\hline
D&\{6,7\}&\{4,5\}&\{10,11\}&\{8,9\}.
\end{array}
\end{equation}

We also record one reusable finish. If $1,2,5,6,7,10$ are selected,
$9,14$ omitted, and exactly two of $15,17,18$ selected, choose $18$ when
selected and otherwise $17$. The three checks are
\[
\begin{array}{c|c|l}
\text{omitted}&v&\text{remaining neighbour checks}\\\hline
15&18&10\text{ retains }7;\ 14,15\text{ each have three selected neighbours}\\
17&18&10,15\text{ retain }7;\ 14\text{ retains }6\text{ active through }1\\
18&17&15\text{ retains }7;\ 14\text{ retains active }6;\ 9\text{ has }2,5,17.
\end{array}
\]
All choices are active nonports.

\paragraph{At most one selected helper.}
If $h=0$, the possible values $k=5,6,7,8$ give $p=0,1,2,3$ and at least
$1,3,5,7$ ready owners, contradicting~\eqref{eq:twenty-ready}.
If $h=1$ and $p=0$, the selected helper is supported. Otherwise the same
count leaves only $k=5,p=1$: square $A$ has three selected owners, omitting
$16$ or $19$, and $B$ has two. Ready owner $13$ forces the omitted witness
to be $5$ or $9$. A selected helper outside $\{0,2\}$ is supported.
For the four remaining possibilities the supported choices are
\[
\begin{array}{c|c|c|l}
D&\text{helper}&v&\text{check}\\\hline
\{5\}&0&13&5\text{ has }0,9,13\\
\{5\}&2&13&5\text{ retains }9\text{ active through }2\\
\{9\}&0&0&4,5\text{ retain }11,13\\
\{9\}&2&2&8\text{ retains }6;\ 9\text{ retains }5\text{ active through }13.
\end{array}
\]

\paragraph{Two selected helpers.}
Equations~\eqref{eq:twenty-budget}--\eqref{eq:twenty-helpers} leave only
$(k,p)=(6,3)$ or $(5,2)$. For $(6,3)$, a full $A$ and two-owner $B$ force
one omission in each of $\{5,9\},\{6,8\},\{4,11\}$.
The single-omission rule would omit helpers $0,2,3$. A full $B$ has four
ready owners and is impossible. Otherwise both squares have three selected
owners, and the ready count forces $A$ to omit $16$ or $19$, with ready
owner $13$. If $B$ has ready owners $14,15$, the three omission pairs are
$\{5,9\},\{6,8\},\{7,10\}$. Helpers $0,1$ are omitted and $2,3$
selected. If $10$ is selected, owner $15$ is supported through that active
mate of omitted $7$. If $10$ is omitted, helper $3$ is supported: $11$
retains $4$, and $10$ retains $7$ active through $15$.
If $B$ has ready owners $17,18$, both $5,9$ are omitted, with one further
omission in $\{7,10\}$, and helper $3$ is omitted. Each of helpers $0,2$
would be supported when selected, since its omitted witness retains its
active ready owner $13$ or $17$. Only helper $1$ could be selected.

For $(k,p)=(5,2)$ the square occupancies are $3+2$. If $A$ omits $12$ or
$13$, ready owners $16,19$ force omissions in $\{6,8\}$ and $\{4,11\}$;
helpers $2,3$ are omitted and $0,1$ selected. Omitted $4$ makes helper $0$
supported through $11$ active at $19$; omitted $11$ makes owner $19$
supported through $4$ active at helper $0$.
If $B$ has ready owners $14,15$, its two omission pairs are
$\{6,8\},\{7,10\}$. Helper $1$ is omitted; helper $0$ would be supported,
so $2,3$ are selected. Omitted $8$ makes helper $2$ supported through $6$
active at $14$; omitted $6$ makes owner $14$ supported through $8$ active
at helper $2$. If $B$ has ready owners $17,18$, helpers $2,3$ are omitted,
$0,1$ selected, and omissions lie in $\{5,9\},\{7,10\}$. Omitted $7$
makes helper $1$ supported through active $10$; omitted $10$ makes owner
$18$ supported through active $7$.

It remains that $A$ omits $16$ or $19$, $B$ has two selected owners and
$13$ is ready. In particular $13$ is selected and active. At least one of
$5,9$ belongs to $D$. If both belong to $D$, each of helpers $0,1,3$ is
supported whenever selected. For helper $0$, selected $4$ retains $11$ and
outside $5$ retains the active owner $13$; helpers $1,3$ have both witness
mates selected. This would leave at most one selected helper. Hence exactly
one of $5,9$ is omitted, and the single-omission rule gives $0\notin C$.

We first exclude the alternative $5\in D$, $9\in C$. Failure to support
owner $13$ forces $9$ inactive, since outside $5$ has just $13,9$ selected.
Thus $2,17\notin C$ and, because $h=2$, helpers $1,3$ are selected.
Write $D=\{5,z\}$, with $z\in\{4,6,7,8,10,11\}$. The following choices
contradict the absence of a supported helper:
\[
\begin{array}{c|c|l}
z&v&\text{checks on its witnesses after deleting }v\\\hline
4,11&1&6,7\text{ retain their selected mates }8,10\\
6,8&3&10,11\text{ retain their selected mates }7,4\\
7&1&6\text{ retains }8;\ \text{outside }7\text{ retains }10
       \text{ active through helper }3\\
10&3&11\text{ retains }4;\ \text{outside }10\text{ retains }7
       \text{ active through helper }1.
\end{array}
\]
The newly outside helper retains an active witness in every row. An omitted
witness with a selected owner has three selected neighbours; otherwise the
displayed active mate verifies its two-neighbour test. This discharges the
entire alternative $5\in D$.

We therefore have $9\in D$ and $5\in C$. Write $D=\{9,z\}$ with the same
six possibilities for $z$. If $2\notin C$, helpers $1,3$ are again selected,
and the preceding table applies unchanged: none of its checks uses the status
of $5$ or $9$. Hence $2\in C$. Outside $9$ retains selected $5$, active
through selected owner $13$, when helper $2$ is deleted. If
$z\notin\{6,8\}$, selected witness $8$ retains selected mate $6$ as well,
so helper $2$ is supported. Only $z=6,8$ remain.

If $z=6$, failure of helper $2$ forces $16\notin C$, since otherwise
selected $8$ retains owner $16$. Helper $3$, if selected, is supported:
its selected witnesses $10,11$ retain selected mates $7,4$.
Thus $3\notin C$, and the two selected helpers are exactly $1,2$.
Deleting helper $1$ is safe: outside $6$ retains selected $8$ active through
helper $2$, and selected $7$ retains selected $10$. A selected owner $14$
would only add a third selected neighbour to outside $6$.
This excludes $z=6$.

If $z=8$, the same check at helper $3$ again gives $3\notin C$, so
$C\cap H=\{1,2\}$. Selected witness $6$ has omitted mate $8$; failure
of helper $1$ therefore forces $14\notin C$, since selected $7$ already
retains $10$. We now have $1,2,5,6,7,10\in C$, $9,14\notin C$, and exactly
two of $15,17,18$ selected. The reusable finish gives a supported owner.

\paragraph{Three selected helpers.}
Let $j$ be the omitted helper; now $p=k-2$. For a ready owner $x+8$, put
$T_x=\{x,s(x),m(x)\}$. The seven sets are
\[
\begin{array}{c|c@{\qquad}c|c}
x&T_x&x&T_x\\\hline
5&\{4,5,9\}&6&\{6,7,8\}\\
7&\{6,7,10\}&8&\{6,8,9\}\\
9&\{5,8,9\}&10&\{7,10,11\}\\
11&\{4,10,11\}&&
\end{array}
\]
Absence of a supported ready owner is equivalent to
\begin{equation}\label{eq:twenty-obstruction}
 U_j(x+8)\subseteq D,\qquad
 U_j(x+8)=\begin{cases}\{m(x)\},&h(x)=j,\\T_x,&h(x)\ne j.\end{cases}
\end{equation}
If the helper is selected, selected $x$ is safe through the helper, and
omitted $x$ is safe if its mate is selected or its helper is active through
$s(x)$. If the helper is omitted, a selected mate is active through its own
selected helper, whereas omission of both witnesses leaves only the owner.
This proves the equivalence, including the pair $\{4,11\}$.

The case $k=8$ has seven ready owners but $p=6$. For a full $A$, the unions
of the sets in~\eqref{eq:twenty-obstruction} have sizes $6,7,6,5$ as
$j=0,1,2,3$; for a full $B$ they have sizes $7,6,6,6$. Thus $B$ cannot be
full. Full $A$ could survive only at $k=7,j=3$, forcing
$D=\{4,5,6,8,9\}$; the three selected owners of $B$ additionally force $7$
omitted. If both squares have three selected owners, $p=4$. Write $A_L$
for omission of $12$ or $13$, $A_R$ for omission of $16$ or $19$, and
analogously $B_L,B_R$. Their forced-union sizes are
\[
\begin{array}{c|rrrr}
&j=0&j=1&j=2&j=3\\\hline
A_L,B_L&8&8&6&6\\
A_L,B_R&7&6&6&6\\
A_R,B_L&6&7&6&5\\
A_R,B_R&5&5&7&7
\end{array}
\]
Each entry is the union of the explicitly displayed sets, and exceeds four.
Therefore $k\le5$; the helper rule gives $k\ge4$.

For $k=4,p=2$, use~\eqref{eq:twenty-h3-p2}. No ready owner's obstruction
fits $D$, so both squares have two selected owners. Failure of the active
helpers gives the following forced omissions and choices:
\[
\begin{array}{c|c|c|l}
j&D&\text{omitted owners}&\text{supported choice}\\\hline
0&\{6,7\}&16,18&\text{as below}\\
1&\{4,5\}&17,19&8\text{ if }14\in C,\text{ otherwise }10\\
2&\{10,11\}&15&6\text{ if }16\in C;\ 7\text{ if }18\in C;\text{ otherwise }9\\
3&\{8,9\}&13,14&5\text{ if }17\in C,\text{ otherwise }7.
\end{array}
\]
For row one, if $17$ is selected, choose the other selected owner of $B$,
either $14$ or $15$. Its outside witness has three selected neighbours,
and its omitted square neighbour retains an active selected witness.
Otherwise $B\cap C=\{14,15\}$. If $A\cap C=\{12,13\}$, choose $13$:
outside $16,19$ each have three selected neighbours. For the other possible
pairs $\{12,19\},\{13,19\}$ choose $5$: outside $0$ retains $4$ active
at $11$, selected $9$ retains helper $2$, and owner $13$ passes the
two-owner rule or retains $19$.
For row two, $8$'s mate $6$ retains selected $14$, and owner $16$ is safe
by the two-owner rule since its same-side opposite $19$ is omitted. If $14$
is omitted, $B\cap C=\{15,18\}$ gives the analogous check for $10$.
For row three, selected $16$ makes $8$ safe for deleting $6$. Otherwise,
selected $18$ gives outside $10$ three selected neighbours for deleting $7$.
If both $16,18$ are omitted, $B\cap C=\{14,17\}$: choice $9$ has outside
helper $2$ retaining $8$ active at $6$, selected $5$ retaining helper $0$,
and selected $17$ retaining $14$. In row four, selected $17$ gives outside
$9$ three selected neighbours for deleting $5$; omitted $13$ is safe by the
two-owner rule. Otherwise $B\cap C=\{15,18\}$ makes $7$ supported.

For $k=5,p=3$, the square occupancies are $3+2$. The obstruction sets show
that a three-owner $A$ must omit $16$ or $19$; its ready owner $13$ forces
$9\in D$ if $j=0$, and $D=\{4,5,9\}$ otherwise. A three-owner $B$ is
possible only at $j=1$, omitting $17$ or $18$, and forces $\{8,10\}\subseteq D$.
The other types force at least four omissions. If $j=2$ or $3$, the first
alternative makes helper $1$ supported, since $6,7,8,10$ are selected.
If $j=1$ and $A$ has three selected owners, choose $6$ unless
$B\cap C=\{14,15\}$: outside helper $1$ retains $7$ active at $10$,
selected $8$ retains helper $2$, and $14$ is safe by the two-owner rule
except in that same-side pair. In that exception choose $14$: $6$ retains
$8$, outside $17$ retains $15$ active at $7$, and outside $18$ has $14,15,10$.
If $j=1$ and $B$ has three selected owners, its selected owners include
$14,15$, and $D=\{8,10,z\}$.
For $z\in\{4,7,11\}$, delete helper $2$: selected $9$ retains selected
$5$, and outside $8$ retains selected $6$, active through selected owner
$14$. If owner $16$ is selected, outside $8$ instead has three selected
neighbours. For $z\in\{5,6,9\}$, delete helper $3$: selected $11$ retains
selected $4$, and outside $10$ retains selected $7$, active through selected
owner $15$. If owner $18$ is selected, outside $10$ has three selected
neighbours. In both groups the newly outside helper retains the first
selected witness, whose displayed mate remains selected.

It remains that $j=0$, $A$ omits $a\in\{16,19\}$, $B$ has two selected
owners, $|D|=3$ and $9\in D$. Owner $13$ is selected and active.
If $a=16$, failure of helper $2$ requires an omission in $\{5,6,8\}$.
The $11$-side of helper $3$ is safe through $19$. Failure of $3$ requires
either both $10,11$ omitted, which would make helper $2$ supported, or $7$
omitted. Thus $D=\{7,9,t\}$ with $t\in\{5,6,8\}$.
For $t=5$ helper $1$ is supported; for $t=6$ witness $5$ is supported.
For $t=8$, failure of helpers $1,3$ forces $14,18$ omitted, so $17$ is
selected and witness $5$ is again supported.
If $a=19$, failure of helper $2$ requires $5$ or $8$ omitted. If $5$ is
omitted but $8$ selected, helper $1$ is supported unless the last omission
is $10$, when helper $3$ is supported. Thus $D=\{8,9,t\}$.
Failure of helper $3$ forces $t\in\{4,7,11\}$, and failure of helper $1$
forces $14$ omitted. For $t=4,11$ the reusable finish applies. For $t=7$,
failure of helper $3$ additionally forces $18$ omitted, so $17$ is selected
and witness $5$ is supported. All these uses of witness $5$ have
$4,5,11,13$ selected, $0,9$ omitted, $13$ with a selected core neighbour,
and either $17$ or $8$ selected. Outside $0$ retains $4$ active at $11$;
outside $9$ either has three selected neighbours, or retains helper $2$
active at $8$. This supplies the required activity checks.

\paragraph{Four selected helpers.}
Write $W=X\cap C$, $w=|W|$. Then $w+k=9$, and the helper rule gives
$p\ge3$, hence $w\le5$, $k\ge4$. A ready owner $x+8$ is safe exactly
when $W\cap T_x\ne\varnothing$. A full $A$ forces $W\subseteq\{7\}$,
so $w+k=9$ would give $w=1,k=8$; but a full $B$ forces
$W\subseteq\{4\}$, a contradiction. Starting with full $B$ gives the same
conclusion. Thus $k\le6$. A three-owner $A$ cannot omit $12$ or $13$,
which would force $w\le2$; it omits $16$ or $19$, forcing
\begin{equation}\label{eq:twenty-WA}
 W\subseteq\{6,7,8,10,11\}.
\end{equation}
A three-owner $B$ must omit $17$ or $18$, forcing
\begin{equation}\label{eq:twenty-WB}
 W\subseteq\{4,5,9,11\}.
\end{equation}
For $k=6$, both apply and $W\subseteq\{11\}$, contradicting $w=3$.
For $k=5$, a three-owner $B$ gives $W=\{4,5,9,11\}$, making helper $0$
supported. Otherwise $W$ is four of the five vertices
in~\eqref{eq:twenty-WA}. If the missing one is $6,7$ or $11$, helper $1$
is supported. For $W=\{6,7,10,11\}$, failure of helper $1$ forces $14$
omitted; witness $7$ is supported, since its helper retains $6$, its mate
$10$ retains helper $3$, and its owner is safe by the two-owner rule.
For $W=\{6,7,8,11\}$, failure of helper $1$ forces $15$ omitted and
witness $6$ passes the corresponding checks.

For $k=4$, both squares have two selected vertices. Indeed a three-owner
$B$ forces $w\le4$, while a three-owner $A$ forces
$W=\{6,7,8,10,11\}$ and makes helper $1$ supported. Exactly three
witnesses are omitted. By the helper rule exactly one helper is isolated;
its two omissions block its neighbouring helpers in the helper cycle, and
the third omission must block the opposite helper. This leaves eight cases.
For each active helper its unique blocking witness is selected, so that
witness's owner must be omitted, except that witness $4$ has no owner edge.
The forced omissions and supported choices are
\[
\begin{array}{c|c|c}
D&\text{omitted owners}&v\\\hline
\{4,5,8\}&19,17,14&7\\
\{4,5,10\}&19,17,15&6\\
\{6,7,11\}&16,18&9\\
\{6,7,9\}&16,18,13&11\\
\{8,9,7\}&14,13,18&5\\
\{8,9,4\}&14,13,19&7\\
\{10,11,6\}&15,16&8\\
\{10,11,5\}&15,17&6
\end{array}
\]
In rows $1,2,4,8$ the selected core pair relevant to the chosen witness is,
respectively,
\[
 \{15,18\},\quad\{14,18\},\quad\{12,19\},\quad\{14,18\}.
\]
The companion and mate of the chosen witness are selected and its owner
retains the other vertex of that pair. In row $3$, witness $9$ has companion
$8$ and mate $5$ selected; owner $17$ is safe by the two-owner rule because
$18$ is omitted. In row $6$, witness $7$ has companion $6$ and mate $10$
selected; owner $15$ is safe because $14$ is omitted.
In row $5$, witness $5$'s helper retains $4$, outside mate $9$ has $2,5,17$
selected since $B\cap C=\{15,17\}$, and omitted $13$ is safe by the
two-owner rule. In row $7$, witness $8$'s helper retains $9$, outside mate
$6$ retains helper $1$ active at $7$, and omitted $16$ is safe by the
two-owner rule. Every choice is an active selected nonport.
All possibilities have contradicted the absence of a supported vertex.
\end{proof}
